% ============================================================
% Relatório final (entrega da disciplina) do projeto
% "Da rede lexical à rede contextual" -- MO438 Redes Complexas, Unicamp.
%
% Compilação:  cd report/entrega-final && latexmk -pdf entrega-final.tex
%
% Versão para leitura do relatório técnico (report/relatorio/), escrita para quem nunca viu o
% repositório. O texto é escrito à mão; figuras e tabelas NÃO são copiadas: vêm do material
% gerado pela etapa `report` em ../relatorio/figuras e ../relatorio/tabelas, e os números do
% corpus e da amostra, de ../relatorio/tabelas/numeros.tex. Duas figuras com os 36 blocos vêm
% da execução separada com todas as camadas (configs/all-layers.yaml), que grava o seu
% material em outputs/experiment/all-layers/relatorio (fora do git). O que faltar aparece
% como "pendente".
% ============================================================
\documentclass[11pt,a4paper]{article}

% ------------------------------------------------------------
% Pacotes
% ------------------------------------------------------------
\usepackage[utf8]{inputenc}
\usepackage[T1]{fontenc}
\usepackage[brazil]{babel}
\usepackage{lmodern}
\usepackage{microtype}
\usepackage{amsmath,amssymb}
\usepackage{graphicx}
\usepackage{booktabs}
\usepackage{tabularx}
\usepackage{array}
\usepackage[table]{xcolor}
\usepackage{enumitem}
\usepackage{geometry}
\usepackage{titlesec}
\usepackage{caption}
\usepackage{tikz}
\usetikzlibrary{arrows.meta,positioning}
% Sem siunitx de propósito: os números gerados já vêm formatados em pt-BR (ex.: 15.104; 0,83).
\usepackage{hyperref}

% ------------------------------------------------------------
% Layout (o do relatório técnico: as figuras geradas foram desenhadas para essa largura)
% ------------------------------------------------------------
\geometry{top=2.2cm, bottom=2.2cm, left=2.0cm, right=2.0cm}

\setlength{\parindent}{0pt}
\setlength{\parskip}{5pt}
\setlength{\emergencystretch}{2.5em}
\linespread{1.04}

\setlist[itemize]{leftmargin=1.35em, itemsep=2pt, topsep=3pt, parsep=0pt}
\setlist[enumerate]{leftmargin=1.55em, itemsep=2pt, topsep=3pt, parsep=0pt}
\setlist[description]{leftmargin=1.2em, itemsep=3pt, topsep=3pt, parsep=0pt}

\titleformat{\section}{\bfseries\large}{\thesection}{0.6em}{}
\titlespacing*{\section}{0pt}{14pt}{5pt}
\titleformat{\subsection}{\bfseries\normalsize}{\thesubsection}{0.6em}{}
\titlespacing*{\subsection}{0pt}{10pt}{3pt}

\captionsetup{font=small, labelfont=bf, margin=0.5cm}
\captionsetup[table]{position=top, skip=4pt}

% Figuras grandes cabem junto do texto; páginas só de figuras começam no alto.
\renewcommand{\topfraction}{0.9}
\renewcommand{\bottomfraction}{0.8}
\renewcommand{\textfraction}{0.08}
\renewcommand{\floatpagefraction}{0.75}
\makeatletter
\setlength{\@fptop}{0pt}
\setlength{\@fpsep}{12pt plus 2fil}
\setlength{\@fpbot}{0pt plus 1fil}
\makeatother

\hypersetup{
    colorlinks=true,
    linkcolor=black,
    citecolor=black,
    urlcolor=blue,
    pdftitle={Da rede lexical à rede contextual: relatório final},
    pdfauthor={Juliana Midlej do Espirito Santo},
}

\newcommand{\repo}{https://github.com/annwith/gender-representation-networks}
\newcommand{\tecnico}{\repo/tree/main/report/relatorio}
\newcommand{\viabilidade}{\repo/blob/main/report/estudo-de-viabilidade/estudo-de-viabilidade.pdf}
\newcommand{\coleta}{https://annwith.github.io/gender-representation-networks/coleta/}

% ------------------------------------------------------------
% Figuras e tabelas não atravessam o início de uma seção (como no relatório técnico)
% ------------------------------------------------------------
\makeatletter
\newcommand{\@flutuantespendentes}{%
    \let\@elt\relax\edef\@tempa{\@botlist\@deferlist\@dbldeferlist}}
\newcommand{\barreiraflutuantes}{\par
    \begingroup\@flutuantespendentes
    \ifx\@tempa\@empty\endgroup\else
        \endgroup\newpage
        \begingroup\@flutuantespendentes
        \ifx\@tempa\@empty\endgroup\else\endgroup\clearpage\fi
    \fi}
\makeatother
\NewCommandCopy{\secaosembarreira}{\section}
\RenewDocumentCommand{\section}{s o m}{%
    \barreiraflutuantes
    \IfBooleanTF{#1}{\secaosembarreira*{#3}}%
        {\IfValueTF{#2}{\secaosembarreira[#2]{#3}}{\secaosembarreira{#3}}}%
    \suppressfloats[t]}

% ------------------------------------------------------------
% Material gerado (pela etapa `report`, na pasta do relatório técnico)
% ------------------------------------------------------------
\newcommand{\gerado}{../relatorio}
% Execução separada com os 36 blocos (configs/all-layers.yaml); fica fora do git.
\newcommand{\geradotodas}{../../outputs/experiment/all-layers/relatorio}
% Números citados no texto: tabelas/numeros.tex define as macros com \newcommand. Os
% \providecommand abaixo só valem enquanto o arquivo não existe ou não define a macro.
% Para citar um número novo, acrescente aqui o \providecommand (a macro já deve existir em
% ../relatorio/relatorio.tex; os testes conferem).
\newcommand{\pendente}{\textit{[pendente]}}
\IfFileExists{\gerado/tabelas/numeros.tex}{\input{\gerado/tabelas/numeros.tex}}{}
\providecommand{\nArtigosCandidatos}{\pendente}
\providecommand{\nArtigos}{\pendente}
\providecommand{\nParagrafos}{\pendente}
\providecommand{\nPalavrasCorpus}{\pendente}
\providecommand{\nTokensCorpus}{\pendente}
\providecommand{\nTiposCorpus}{\pendente}
\providecommand{\nVertices}{\pendente}
\providecommand{\nTipos}{\pendente}
\providecommand{\nSequencias}{\pendente}
\providecommand{\nTokensProcessados}{\pendente}
\providecommand{\nParagrafosNucleo}{\pendente}
\providecommand{\nVerticesNucleo}{\pendente}
\providecommand{\nVerticesAlvo}{\pendente}
\providecommand{\nVerticesControle}{\pendente}
\providecommand{\nVerticesMultitema}{\pendente}
\providecommand{\nAlvosMantidos}{\pendente}
\providecommand{\nGruposPrefixo}{\pendente}
\providecommand{\nTiposVocab}{\pendente}
\providecommand{\nTiposVocabCorpus}{\pendente}
\providecommand{\difMaxPrefixo}{\pendente}
\providecommand{\concordanciaPredicao}{\pendente}
\providecommand{\dataExecucao}{\pendente}

\newsavebox{\caixafigura}
% Inclui um PDF no tamanho natural, reduzindo só se for mais largo que a coluna.
\newcommand{\incluirlimitado}[1]{%
    \sbox{\caixafigura}{\includegraphics{#1}}%
    \ifdim\wd\caixafigura>\linewidth
        \includegraphics[width=\linewidth]{#1}%
    \else
        \usebox{\caixafigura}%
    \fi}
% Ocupa o lugar de um material ainda não gerado (só aparece se o arquivo faltar).
\newcommand{\caixapendente}[1]{%
    \fcolorbox{black!45}{black!4}{%
        \parbox{0.9\linewidth}{\centering\small\color{black!75}%
            \rule{0pt}{1.8em}#1\par\vspace{1.1em}}}}
% \figuraopcional{arquivo sem extensão}{legenda}{rótulo}: arquivo em ../relatorio/figuras/
\newcommand{\figuraopcional}[3]{%
    \begin{figure}[htbp]
        \centering
        \IfFileExists{\gerado/figuras/#1.pdf}{%
            \incluirlimitado{\gerado/figuras/#1.pdf}%
        }{%
            \caixapendente{\textbf{Figura pendente:}
                \texttt{\detokenize{figuras/#1.pdf}} (etapa \texttt{report})}%
        }
        \caption{#2}\label{#3}
    \end{figure}}
% \figuratodas{arquivo sem extensão}{legenda}{rótulo}: a mesma figura, desenhada pela execução
% com os 36 blocos (outputs/experiment/all-layers/relatorio/figuras/)
\newcommand{\figuratodas}[3]{%
    \begin{figure}[htbp]
        \centering
        \IfFileExists{\geradotodas/figuras/#1.pdf}{%
            \incluirlimitado{\geradotodas/figuras/#1.pdf}%
        }{%
            \caixapendente{\textbf{Figura pendente:}
                \texttt{\detokenize{figuras/#1.pdf}} da execução com todos os blocos
                (\texttt{configs/all-layers.yaml}, etapa \texttt{report})}%
        }
        \caption{#2}\label{#3}
    \end{figure}}
% \tabelaopcional{arquivo sem extensão}{legenda}{rótulo}: arquivo em ../relatorio/tabelas/
\newcommand{\tabelaopcional}[3]{%
    \begin{table}[htbp]
        \centering
        \caption{#2}\label{#3}
        \IfFileExists{\gerado/tabelas/#1.tex}{%
            \small\input{\gerado/tabelas/#1.tex}%
        }{%
            \caixapendente{\textbf{Tabela pendente:}
                \texttt{\detokenize{tabelas/#1.tex}} (etapa \texttt{report})}%
        }
    \end{table}}

% ------------------------------------------------------------
% Notação
% ------------------------------------------------------------
\newcommand{\tok}[1]{\texttt{#1}}              % um token, como aparece no texto
\newcommand{\esp}{\textvisiblespace}           % espaço inicial de um token
\newcommand{\cod}[1]{\texttt{#1}}              % código, arquivos, comandos
\newcommand{\N}{\mathcal{N}}

% ============================================================
\begin{document}

\begin{center}
    {\large\bfseries MO438 -- Redes Complexas: Relatório Final}\\[6pt]
    {\LARGE\bfseries Da rede lexical à rede contextual}\\[4pt]
    {\large Como um Transformer reorganiza as ocorrências de tokens, visto por redes
    complexas}\\[10pt]
    Juliana Midlej do Espirito Santo (RA 200208)\\[2pt]
    {\small Unicamp. Números da execução de \dataExecucao{}.}
\end{center}

\begin{abstract}
\noindent
Na entrada de um Transformer, todas as ocorrências de um mesmo token são representadas pelo
mesmo vetor; ao atravessar as camadas, cada ocorrência passa a carregar o texto em que
aparece. Este trabalho descreve essa passagem com redes complexas. Os vértices são
\nVertices{} ocorrências de tokens em parágrafos da Wikipédia em português, de oito temas, e
cada vértice se liga aos 10 vértices de vetor mais parecido (k vizinhos mais próximos pelo
cosseno). Os mesmos vértices recebem quatro vetores do modelo Qwen3-4B-Base: o
\emph{embedding} de entrada e as saídas dos blocos 1, 18 e 36. As quatro redes têm, portanto,
os mesmos vértices e podem ser comparadas vértice a vértice. No bloco 1, a rede ainda é
praticamente a rede lexical: as ocorrências de um token continuam vizinhas umas das outras.
A reorganização vem depois e é grande. No bloco 36, uma ocorrência mantém em média menos de 2
dos seus 10 vizinhos lexicais, a fração de vizinhos do mesmo token cai de 0,70 para 0,25
(ainda muito acima do acaso), surgem \emph{hubs} e as distâncias encurtam. As comunidades
deixam de seguir o token e passam a seguir o contexto: nos blocos do meio (do 12 ao 19,
aproximadamente), elas concordam mais com o artigo, o parágrafo e o tema de origem do que com
o próprio token; nos últimos blocos, essa concordância cai e cresce a concordância com o
próximo token. Palavras ambíguas, como \emph{banco}, \emph{órgão} e \emph{matriz}, separam as
suas ocorrências por tema e por sentido mais do que palavras de sentido estável. O texto
descreve também a coleta dos dados, a construção das redes, o tratamento dos empates e as
limitações do estudo.
\end{abstract}

% ============================================================
\section{Introdução}\label{sec:introducao}
% ============================================================

\previsto{Cerca de 1 página. Motivação (na entrada de um Transformer, todas as ocorrências de
um token têm o mesmo vetor; ao longo das camadas elas incorporam o contexto). A ideia do
projeto: uma sequência de redes $k$-NN sobre \emph{os mesmos vértices} (ocorrências de tokens),
uma por representação, comparadas vértice a vértice. As perguntas P1--P4 com as hipóteses
H1--H4, em uma frase cada (relatório técnico, Seção 2). Contribuições e organização do texto.}

% ============================================================
\section{Dados: da Wikipédia à amostra de ocorrências}\label{sec:dados}
% ============================================================

Os vértices das redes são ocorrências de tokens em parágrafos reais da Wikipédia em
português. Esta seção descreve o caminho dos dados em duas etapas. Primeiro, montamos um
\textbf{corpus}: \nArtigos{} artigos de oito temas, limpos e divididos em \nParagrafos{}
parágrafos de prosa corrida. Depois, sorteamos desse corpus uma \textbf{amostra} de
\nVertices{} ocorrências, que são os vértices de todas as redes. A página interativa
\url{\coleta} mostra a coleta, a amostra e a construção das redes passo a passo.

% ------------------------------------------------------------
\subsection{Temas e categorias da Wikipédia}\label{sec:temas}

A Wikipédia organiza os artigos em \textbf{categorias}. Uma categoria contém artigos e outras
categorias (subcategorias), e essas ligações formam uma hierarquia: a categoria Física contém
subcategorias como Cosmologia e Efeitos físicos, que têm as suas próprias subcategorias, e
assim por diante. Usamos essa hierarquia
para definir \textbf{temas} de forma reproduzível: um tema é o conjunto de artigos alcançados
a partir de uma categoria raiz.

São oito temas (Tabela~\ref{tab:temas}). Eles foram escolhidos distantes entre si, para que o
vocabulário de cada um seja característico, e de modo que as palavras ambíguas da P4 tenham
sentidos diferentes em pelo menos dois deles: \emph{banco} é instituição financeira em
Economia, banco de dados em Computação e banco de areia em Geografia. O tema tem dois papéis
na análise: é um dos rótulos de contexto da P3 e serve como indicador aproximado do sentido de
uma palavra ambígua na P4. O modelo nunca vê o tema; ele é só um rótulo usado na análise.

\begin{table}[htbp]
\centering
\caption{Os oito temas. ``Artigos'' conta os artigos mantidos no corpus, vindos da árvore da
categoria raiz e das categorias extras (Seção~\ref{sec:extras}). A última coluna resume de onde
vieram os 150 artigos da raiz, pelas subcategorias que mais contribuíram.}\label{tab:temas}
\small
\begin{tabularx}{\linewidth}{@{}l l r r >{\raggedright\arraybackslash}X@{}}
\toprule
\textbf{Tema} & \textbf{Categoria raiz} & \multicolumn{2}{c}{\textbf{Artigos}} &
\textbf{O que a busca trouxe da raiz} \\
\cmidrule(lr){3-4}
& & raiz & extras & \\
\midrule
Física & Física & 150 & -- & conceitos e efeitos físicos, astronomia, história e
organizações da física; também prêmios (33 artigos) e livros de física \\
Biologia & Biologia & 150 & 240 & membros diretos da raiz (78), biotecnologia (36),
biomedicina, divisões da biologia \\
Economia & Economia & 150 & 300 & membros diretos da raiz (34), terminologia e indicadores
econômicos, economia financeira, teorias e modelos \\
Política & Política & 150 & -- & membros diretos da raiz (45), ciência política (37),
direitos humanos (29), constituições, corrupção \\
Computação & Ciência da computação & 150 & 173 & membros diretos da raiz (123);
armazenamento de dados, cibernética \\
Música & Música & 150 & 60 & armazenamento de áudio, gêneros, subculturas e instrumentos
musicais, história e indústria da música \\
Esporte & Futebol & 150 & -- & terminologia, regras e história do futebol, futebol por
continente e por país, jogos eletrônicos e hinos de futebol \\
Geografia & Geografia & 150 & 180 & cartografia, geomática, acidentes geográficos, termos
geográficos \\
\bottomrule
\end{tabularx}
\end{table}

% ------------------------------------------------------------
\subsection{Busca pelas categorias}\label{sec:busca}

Para cada tema, uma \textbf{busca em largura} parte da categoria raiz e desce pelas
subcategorias até a \textbf{profundidade 2}: a raiz é a profundidade 0, as suas subcategorias
são a profundidade 1 e as subcategorias delas, a 2. A busca guarda os títulos de artigos que
encontra, até \textbf{600 títulos por tema}, e segue três regras:
\begin{itemize}
    \item só entram páginas do domínio principal (artigos), e não páginas de usuário,
    predefinições ou discussões;
    \item não abre subcategorias de manutenção, cujo nome começa com ``!'' na Wikipédia em
    português (esboços, avaliação de qualidade, predefinições, portais), nem subcategorias de
    listas (``Listas de\ldots'');
    \item descarta os títulos ``Lista de\ldots'', que não são prosa.
\end{itemize}
A busca registra também, para cada título, \emph{todos} os temas que chegam a ele, o que
permite descartar as páginas que pertencem a mais de um tema (Seção~\ref{sec:descartes}).
Depois, os títulos de cada tema são embaralhados com uma semente fixa e baixados em lotes de
50, e os artigos que passam pelos filtros entram até completar \textbf{150 por tema}. Todas as
respostas da API ficam guardadas em disco, e a revisão exata de cada artigo (\cod{revid}) fica
registrada; assim o mesmo texto pode ser obtido de novo, mesmo que os artigos mudem.

\paragraph{Por que limitar a busca, e o que isso custa.} A árvore de categorias da Wikipédia
não é uma taxonomia limpa: em poucos níveis, uma categoria de ciência chega a pessoas, lugares
e obras. Por isso a busca vai só até a profundidade 2 e para em 600 títulos. Esse corte, porém,
\textbf{não é aleatório}. A API lista os membros de uma categoria em ordem alfabética, e a busca
abre as subcategorias nessa ordem; quando o limite é atingido, as que ainda estão na fila nunca
são abertas. Os oito temas atingiram o limite, e a busca abriu só uma parte das subcategorias
que encontrou: de 7 de 79 (Computação) a 98 de 439 (Música). Em Biologia, Política e
Computação, o corte veio antes de terminar a profundidade 1. Em Futebol, a busca parou no meio
da lista de subcategorias por país, em ``Futebol das Bahamas'': dos 150 artigos mantidos, 29
vêm de países de Afeganistão a Azerbaijão e nenhum do Brasil ou de Portugal. A profundidade
também varia: em Computação, 123 dos 150 artigos são membros diretos da raiz; em Geografia,
105 vêm da profundidade 2. A última coluna da Tabela~\ref{tab:temas} mostra o efeito: cada tema
representa as subcategorias mais próximas da raiz e alfabeticamente primeiras, e não a área
inteira, e inclui artigos pouco típicos (prêmios em Física; em Futebol, canções pop que
viraram hinos de torneios, como ``Colors'', de Jason Derulo e Maluma). Como o objetivo é ter temas bem separados, e não medir o vocabulário de cada
área, isso é aceitável, mas os nomes dos temas devem ser lidos como ``artigos próximos da
categoria raiz''.

% ------------------------------------------------------------
\subsection{Categorias extras}\label{sec:extras}

Alguns sentidos das palavras ambíguas quase não aparecem no tema que os representa. Num estudo
preliminar (Seção~\ref{sec:amostra}), \emph{nota} no sentido de cédula tinha só 7 ocorrências
em Economia, \emph{carga} tributária 6, \emph{órgão} do corpo 6 em Biologia e \emph{banco} de
areia 4 em Geografia. Para esses casos, o corpus inclui \textbf{16 categorias extras}, cada uma
ligada a um dos oito temas (Tabela~\ref{tab:extras}). A busca nelas vai só até a profundidade
1 e fica com até 60 artigos por categoria. Um artigo de uma categoria extra pertence, para
todos os efeitos, ao tema a que ela está ligada: um artigo de Numismática é um artigo de
Economia. Por isso o número de artigos varia entre os temas (de 150 em Física, Política e
Esporte a 450 em Economia).

\begin{table}[htbp]
\centering
\caption{As 16 categorias extras e o sentido que cada uma reforça.}\label{tab:extras}
\small
\begin{tabularx}{\linewidth}{@{}l >{\raggedright\arraybackslash}p{6.2cm} X@{}}
\toprule
\textbf{Tema} & \textbf{Categoria extra} & \textbf{Sentido reforçado} \\
\midrule
Economia & Numismática & \emph{nota} (cédula) \\
 & Impostos; Direito tributário & \emph{carga} (tributária) \\
 & Sistema bancário & \emph{banco} (instituição financeira) \\
 & Indústria & \emph{planta} (industrial) \\
\midrule
Biologia & Anatomia humana; Órgãos (anatomia) & \emph{órgão} (do corpo) \\
 & Botânica & \emph{planta} (vegetal) \\
 & Taxonomia & \emph{família} (grupo taxonômico) \\
\midrule
Computação & Bancos de dados & \emph{banco} (de dados) \\
 & Redes de computadores & \emph{rede} (de computadores) \\
 & Matrizes & \emph{matriz} (matemática) \\
\midrule
Geografia & Geomorfologia; Oceanografia & \emph{banco} (de areia), \emph{massa} (de água) \\
 & Hidrografia & \emph{rede} (hidrográfica) \\
\midrule
Música & Teoria musical & \emph{nota} (musical) \\
\bottomrule
\end{tabularx}
\end{table}

% ------------------------------------------------------------
\subsection{Páginas descartadas: multitema, biografias e outras}\label{sec:descartes}

Um artigo encontrado pela busca é descartado em cinco casos (Tabela~\ref{tab:corpus-temas}):
\begin{itemize}
    \item \textbf{Página multitema}: o artigo é alcançado a partir de mais de um tema,
    contando as categorias extras. Exemplos reais: ``Prêmio Nobel'' (Física e Economia),
    ``Harmonices Mundi'' (Física e Música) e ``Big science'' (Física e Política). Uma página
    assim mistura o vocabulário de dois temas e borraria justamente o rótulo de tema usado na
    P3 e na P4. Foram 108 descartes.
    \item \textbf{Biografia}, reconhecida pela caixa de informações do artigo (predefinições
    \cod{Info/Biografia}, \cod{Info/Futebolista}, \cod{Info/Político}, \cod{Info/Músico},
    \cod{Info/Cientista} e outras 9). Biografias aparecem em quase todas as árvores de
    categorias e falam de datas, lugares e carreiras, um vocabulário que não é o do tema.
    Foram 68.
    \item \textbf{Sem parágrafos}: depois da limpeza (abaixo), não sobra nenhum parágrafo de
    prosa, em geral porque o artigo é um esboço curto ou é feito de tabelas e listas. Foram
    295.
    \item \textbf{Desambiguação} (4) e \textbf{duplicada} (3, o mesmo artigo por dois
    caminhos).
\end{itemize}
Dos \nArtigosCandidatos{} candidatos examinados, ficaram \nArtigos{}.

% ------------------------------------------------------------
\subsection{Limpeza dos parágrafos}\label{sec:limpeza}

Os artigos chegam em \emph{wikitext}, a linguagem de marcação da Wikipédia, que mistura o texto
com referências, tabelas, fórmulas, caixas de informação e predefinições. A limpeza tem
quatro passos.

\textbf{1. Cortar as seções de apoio.} As seções Referências, Ligações externas, Ver também,
Bibliografia, Notas, Fontes, Leitura adicional (e variantes) são removidas com as suas
subseções. Seções de conteúdo que vêm depois delas são mantidas.

\textbf{2. Converter em texto, removendo o que não é prosa.} A conversão usa a biblioteca
\cod{mwparserfromhell}, mas a conversão padrão dela mantém o conteúdo de marcas que não são
texto corrido. Antes de converter, a árvore de marcação é reescrita em três pontos:
\begin{itemize}
    \item referências (\cod{<ref>}), galerias, tabelas, blocos pré-formatados e links para arquivos
    e categorias são removidos com todo o conteúdo. Sem isso, o texto das referências entrava
    no meio das frases, colado à palavra anterior (``A energiaSilva, Livro, 2001 é\ldots''),
    defeito que encontramos no estudo preliminar;
    \item linhas com fórmulas (\cod{<math>}, \cod{<chem>}) e itens de lista são descartados
    inteiros, porque apagar só a fórmula deixaria um buraco na frase (``a energia é dada por ,
    onde\ldots''). Fórmulas triviais, de um símbolo ou um número, são mantidas como texto;
    \item predefinições cujo texto visível é um argumento (números formatados, conversões de
    unidades) são trocadas por esse texto; as demais (caixas de informação, navegação,
    pronúncia) são descartadas, e os restos que deixam, como ``( )'' e ``, ,'', são limpos.
\end{itemize}

\textbf{3. Ficar só com prosa corrida.} Um parágrafo é mantido se tem pelo menos 40 palavras,
termina como frase (``.'', ``!'' ou ``?'', talvez seguidos de aspas ou parênteses), não começa
com marcação de tabela ou lista e tem no máximo 8\% de dígitos (o que elimina tabelas de
resultados e cronologias). Parágrafos repetidos, comuns porque há texto copiado entre artigos,
são removidos comparando os 80 primeiros caracteres.

\textbf{4. Dividir em sentenças.} Cada parágrafo é dividido em sentenças por expressão
regular, com uma lista de abreviações do português (``Sr.'', ``séc.'', ``p.ex.'', iniciais de
nomes), para dar o rótulo de sentença usado na P3.

O corpus final tem \nArtigos{} artigos, \nParagrafos{} parágrafos, \nPalavrasCorpus{}
palavras e \nTokensCorpus{} tokens do modelo, de \nTiposCorpus{} tokens distintos
(Tabela~\ref{tab:corpus-temas}).

\tabelaopcional{corpus_temas}
    {O corpus por tema: artigos, parágrafos e palavras mantidos, e artigos descartados por
    motivo. Os artigos incluem os das categorias extras.}
    {tab:corpus-temas}

% ------------------------------------------------------------
\subsection{Tokens, tipos, ocorrências e categorias de token}\label{sec:tokens}

O modelo não lê palavras, e sim \textbf{tokens}, as unidades produzidas pelo seu
\emph{tokenizer}. Um token pode ser uma palavra inteira, um pedaço de palavra, um sinal de
pontuação ou um dígito. No Qwen3, o espaço antes de uma palavra faz parte do token seguinte
(\tok{\esp banco} é um token só; o símbolo \esp{} marca esse espaço) e cada dígito é um token
separado. Usamos três termos:
\begin{itemize}
    \item \textbf{tipo de token}: um elemento do vocabulário do modelo, como \tok{\esp banco};
    o número de vezes que um tipo $t$ aparece na amostra é a sua frequência $f_t$;
    \item \textbf{ocorrência}: uma aparição de um tipo numa posição de um parágrafo. Cada
    ocorrência da amostra é um \textbf{vértice};
    \item \textbf{palavra}: uma sequência de letras ou dígitos no texto original, antes da
    tokenização. Uma palavra pode virar vários tokens.
\end{itemize}

Cada ocorrência recebe uma \textbf{categoria de token}, conforme a sua relação com a palavra a
que pertence (Tabela~\ref{tab:exemplos-tokens}):
\begin{description}
    \item[palavra inteira] o token cobre sozinho uma palavra (\tok{\esp banco},
    \tok{\esp de});
    \item[início de palavra] o primeiro token de uma palavra dividida em vários
    (\tok{\esp torn} em ``tornou'');
    \item[continuação] um token que continua uma palavra já começada (\tok{ou} em
    ``tornou'');
    \item[pontuação] um token sem letras nem dígitos, mas com algum caractere visível
    (\tok{,}, \tok{\esp(});
    \item[número] um token só de dígitos; ``1990'' vira quatro tokens de número.
\end{description}
Tokens formados só por espaço não viram vértices: eles não têm conteúdo próprio e
apareceriam em quase todo parágrafo.

\begin{table}[htbp]
\centering
\caption{Exemplos reais do \emph{tokenizer} do Qwen3-4B-Base e as categorias atribuídas.
O símbolo \esp{} marca o espaço que faz parte do token.}\label{tab:exemplos-tokens}
\small
\begin{tabular}{@{}l l l l@{}}
\toprule
\textbf{Texto} & \textbf{Tokens} & \textbf{Categoria} & \textbf{Palavra} \\
\midrule
``\esp 1990'' & \tok{\esp} & só espaço (não é vértice) & -- \\
 & \tok{1}, \tok{9}, \tok{9}, \tok{0} & número & 1990 \\
\midrule
``\esp tornou-se'' & \tok{\esp torn} & início de palavra & tornou \\
 & \tok{ou} & continuação & tornou \\
 & \tok{-se} & palavra inteira & se \\
\midrule
``\esp(banco)'' & \tok{\esp(} & pontuação & -- \\
 & \tok{ban} & início de palavra & banco \\
 & \tok{co} & continuação & banco \\
 & \tok{)} & pontuação & -- \\
\bottomrule
\end{tabular}
\end{table}

O terceiro exemplo mostra um cuidado necessário na P4: ``banco'' entre parênteses vira
\tok{ban}${}+{}$\tok{co}, e ``Banco'' com maiúscula é outro tipo (\tok{\esp Banco}). Por
isso as palavras ambíguas são contadas \emph{só} no tipo minúsculo precedido de espaço, para
que todas as ocorrências comparadas comecem do mesmo vetor de entrada.

O tokenizer do Qwen3 foi treinado sobretudo em inglês, chinês e código, e divide muito o
português: no estudo preliminar, foram 1,78 tokens por palavra, e só 28\% dos tokens de um
texto corrido eram palavras inteiras. Por isso todos os resultados são separados por
categoria de token, e as interpretações sobre sentido se limitam a palavras inteiras.

% ------------------------------------------------------------
\subsection{A amostra híbrida}\label{sec:amostra}

\paragraph{Por que não usar só parágrafos inteiros.} Antes de montar o experimento, um
\href{\viabilidade}{estudo de viabilidade} simulou a amostragem num corpus preliminar. Ele
mostrou que reunir cerca de 15 mil vértices é fácil (bastam de 100 a 160 parágrafos
inteiros), mas que parágrafos inteiros não servem para a P4: pela lei de Zipf, quase nenhuma
palavra de conteúdo se repete o bastante. Com 15 mil vértices tirados de parágrafos inteiros,
só de 4 a 9 palavras de conteúdo apareciam pelo menos 10 vezes e em 3 temas, e nenhuma palavra
ambígua chegava a 10 ocorrências em 2 temas. A solução foi uma \textbf{amostra híbrida}: um
núcleo de parágrafos inteiros, para as perguntas sobre a rede como um todo, e três estratos de
palavras escolhidas, sorteadas em outros parágrafos, para a P4.

A amostra limita cada tipo a $f_{\max}=50$ ocorrências e usa uma semente fixa. Tem quatro
estratos:

\paragraph{1. Núcleo (P1 a P3).} Para cada tema, sorteamos 5 artigos e 2 parágrafos de cada
um: 80 parágrafos. Todos os tokens desses parágrafos (menos os de só espaço) são vértices,
inclusive pedaços de palavra, números e pontuação. Dois parágrafos por artigo fazem ``artigo''
e ``parágrafo'' serem rótulos diferentes na P3. O limite de 50 ocorrências por tipo é
aplicado por sorteio uniforme depois de juntar todo o núcleo: dos 11.935 tokens dos 80
parágrafos, saíram 1.915 de 25 tipos muito frequentes (a vírgula passou de 483 para 50
ocorrências; \tok{\esp de}, de 391 para 50). Assim as ocorrências mantidas de ``de'' se
espalham por todos os parágrafos, em vez de serem as dos primeiros. O núcleo tem
\nVerticesNucleo{} vértices e é o único estrato em que o parágrafo aparece inteiro; por isso a
P3 é analisada só nele.

\paragraph{2. Palavras-alvo (P4).} Doze palavras com sentidos diferentes em temas diferentes,
cada uma com uma \textbf{cota por tema de sentido}: 2 ou 3 temas, com 16 a 25 ocorrências em
cada (Tabela~\ref{tab:cotas}). As ocorrências são sorteadas fora do núcleo, com no máximo 1 por
parágrafo e 3 por artigo, para que um único artigo não domine uma palavra. Uma palavra que não
chega a 10 ocorrências em pelo menos 2 temas sai da análise: foi o caso de \emph{planta}, que
teve só 5 ocorrências em Economia. Ficaram \nAlvosMantidos{} alvos, com \nVerticesAlvo{}
vértices sorteados. Algumas cotas não foram cumpridas: \emph{banco} teve 7 ocorrências em
Geografia, \emph{órgão} 7 em Política e \emph{família} 11 em Política.

\paragraph{3. Controles (P4).} Seis palavras de sentido estável: \emph{ano}, \emph{século},
\emph{grande}, \emph{primeiro}, \emph{importante} e \emph{água}. Elas têm a \textbf{mesma
alocação} dos alvos (duas com três temas e cotas 17/17/16, quatro com dois temas e cotas
25/25), para que uma diferença entre alvos e controles não venha do desenho da amostra. São
\nVerticesControle{} vértices.

\paragraph{4. Palavras multitema (P4).} 150 palavras de conteúdo que aparecem pelo menos 5
vezes em pelo menos 3 temas fora do núcleo, cada uma completada até 20 ocorrências, sorteadas
em rodízio entre os temas. Exemplos: \emph{metal}, \emph{luz}, \emph{guerra},
\emph{material}, \emph{forte}, \emph{motor}, \emph{banda}, \emph{plano}. Elas estendem a P4
para além das 18 palavras escolhidas à mão, com palavras comuns a vários temas que não foram
escolhidas por serem ambíguas. São \nVerticesMultitema{} vértices. (``Multitema'' tem aqui um
sentido diferente do da Seção~\ref{sec:descartes}: lá, é uma \emph{página} que pertence a
mais de um tema e foi descartada; aqui, é uma \emph{palavra} que aparece em vários temas.)

\begin{table}[htbp]
\centering
\caption{Cotas por tema de sentido das palavras-alvo e de controle, como definidas no desenho
(ocorrências por tema).}\label{tab:cotas}
\small
\begin{tabular}{@{}l l@{\hspace{2.5em}}l l@{}}
\toprule
\textbf{Alvo} & \textbf{Cotas} & \textbf{Controle} & \textbf{Cotas} \\
\midrule
banco & Economia 17, Computação 17, Geografia 16 & ano & Economia 17, Esporte 17, Música 16 \\
campo & Física 17, Esporte 17, Geografia 16 & século & Física 17, Geografia 17, Política 16 \\
órgão & Biologia 17, Música 17, Política 16 & grande & Física 25, Biologia 25 \\
estado & Física 25, Política 25 & primeiro & Esporte 25, Computação 25 \\
carga & Física 25, Economia 25 & importante & Biologia 25, Economia 25 \\
rede & Computação 25, Geografia 25 & água & Biologia 25, Geografia 25 \\
nota & Música 25, Economia 25 & & \\
capital & Economia 25, Política 25 & & \\
família & Biologia 25, Política 25 & & \\
massa & Física 25, Geografia 25 & & \\
matriz & Biologia 25, Computação 25 & & \\
planta & Biologia 25, Economia 25 & & \\
\bottomrule
\end{tabular}
\end{table}

Nos três estratos de palavras escolhidas, o parágrafo da ocorrência também passa pelo modelo,
mas só como contexto: o único vértice dele é a ocorrência sorteada. Ao todo, a amostra tem
\nVertices{} vértices de \nTipos{} tipos, em \nSequencias{} parágrafos de 1.172 artigos
(Figura~\ref{fig:amostra-composicao}). Como o núcleo tem três quartos dos vértices, as medidas
globais descrevem sobretudo texto corrido. Os estratos acrescentam palavras inteiras, e as
palavras inteiras são 48\% dos vértices, contra 28\% num texto corrido.

\figuraopcional{amostra_composicao}
    {Composição da amostra: vértices por estrato, por tema, por faixa de frequência do tipo na
    amostra ($f_t$) e por categoria de token.}
    {fig:amostra-composicao}

% ------------------------------------------------------------
\subsection{Anotação dos sentidos}\label{sec:dados-sentidos}

O tema é só um indicador aproximado do sentido: um artigo de Computação pode citar o Banco
Central. Para medir a separação por sentido de verdade, as 508 ocorrências das 11 palavras-alvo
(inclusive as que caíram no núcleo) foram \textbf{anotadas pelo modelo de linguagem Claude
Opus~5.5}, sem revisão humana. Cada ocorrência foi lida com o trecho do parágrafo em volta, sem
o tema e sem o título do artigo, para que o rótulo registrasse o sentido e não o tema. Os
sentidos de cada palavra (de 3 a 7) foram definidos durante a leitura; por exemplo,
\emph{órgão} tem os sentidos ``órgão biológico'', ``instrumento musical'' e ``órgão
institucional''. Como o modelo analisado só vê o texto à esquerda da palavra, cada ocorrência
recebeu também uma marca que diz se esse trecho à esquerda já decide o sentido (442 das 508).
A separação por sentido é medida nos dois recortes. A anotação e o protocolo estão em
\cod{data/annotations/} no repositório.

% ============================================================
\section{Modelo e construção das redes}\label{sec:redes}
% ============================================================

\previsto{Cerca de 1,5 página, escrita para quem nunca viu o projeto (a parte que mais
confundiu na entrega parcial). Qwen3-4B-Base e o que é cada representação: \emph{embedding} de
entrada (\cod{lex}) e saídas dos blocos 1, 18 e 36. Categorias de token (palavra inteira,
início, continuação, pontuação, número). Construção: cosseno, $k=10$ vizinhos, rede
direcionada e simetrizações por união e mútua, empates exatos da rede lexical (ocorrências do
mesmo token têm cosseno 1) e o desempate com sementes. A rede lexical de tipos sobre o
vocabulário inteiro como referência. Um diagrama pequeno da construção e o link para a página
interativa (\url{\coleta}, aba de construção). Relatório técnico, Seções 3, 5 e 6.}

% ============================================================
\section{Medidas}\label{sec:metodos}
% ============================================================

Esta seção define as medidas usadas nos resultados e diz a que pergunta cada uma responde.

% ------------------------------------------------------------
\subsection{Métricas estruturais}\label{sec:met-estruturais}

Em cada rede calculamos as métricas usuais, com o igraph~\cite{csardi2006}: número de vértices
$n$ e de arestas $m$, grau médio $\langle k\rangle=2m/n$, densidade, reciprocidade (a fração
dos arcos $i\to j$ para os quais $j\to i$ também existe), clusterização global (a
transitividade, $3\times$ triângulos sobre trios conexos) e local média
($C_i=2e_i/[k_i(k_i-1)]$, com $e_i$ o número de arestas entre os vizinhos de $i$), número de
componentes conexos e fração de vértices no maior, distância média e diâmetro, por busca em
largura exata a partir de todos os vértices do maior componente. As métricas não direcionadas
usam a rede por união.

Duas observações ajudam a ler essas métricas numa rede k-NN. Primeiro, o número de arestas da
rede por união fica entre $nk/2$ (todas as escolhas recíprocas) e $nk$ (nenhuma recíproca),
então o grau médio fica entre $k$ e $2k$ e a densidade é baixíssima por construção; o que
interessa é como esses valores variam entre as camadas. Segundo, como todo vértice tem grau de
saída $k$, a informação sobre \emph{hubs} está no \textbf{grau de entrada}. Reportamos a
distribuição dele (a probabilidade $P(K\ge d)$ de um vértice ter grau de entrada pelo menos
$d$) e a sua assimetria (\emph{skewness}), a medida usual de
\emph{hubness}~\cite{radovanovic2010}: assimetria alta significa que poucos vértices são
escolhidos por muitos.

% ------------------------------------------------------------
\subsection{Troca de vizinhos e piso de ruído (P1 e P2)}\label{sec:met-jaccard}

A vizinhança de uma ocorrência $i$ nas redes $a$ e $b$ (por exemplo, \cod{lex} e \cod{L18})
é comparada pelo índice de Jaccard,
\begin{equation}
    J_i(a,b)=\frac{|\N_i^{(a)}\cap\N_i^{(b)}|}{|\N_i^{(a)}\cup\N_i^{(b)}|}.
\end{equation}
Com $k=10$ vizinhos nas duas redes e $c$ vizinhos em comum, $J_i=c/(20-c)$: $J=1$ quando nada
muda, $J=0{,}11$ quando sobram 2 vizinhos e $J=0$ quando nenhum se mantém. A
\textbf{mudança} é $M_i=1-J_i$.

Na rede lexical, parte dos vizinhos é sorteada entre empatados, e duas sementes diferentes já
dão $J<1$ sem nenhuma mudança de representação. O \textbf{piso de ruído} $J^{\text{ruído}}_i$
é o Jaccard médio entre redes lexicais de sementes diferentes (10 pares). Ele depende muito da
frequência: para um tipo com 50 ocorrências, cada sorteio escolhe 10 de 49 cópias idênticas, e
duas escolhas independentes têm em média só 2 em comum ($J^{\text{ruído}}=0{,}12$); para tipos
de 10 a 49 ocorrências, $0{,}39$; para tipos raros, cerca de $0{,}75$. A \textbf{mudança acima
do piso} é $M^{+}_i=J^{\text{ruído}}_i-J_i(\mathrm{lex},\ell)$: a parte da troca de vizinhos que
o sorteio não explica. Como complemento, o Jaccard por composição $J^w$ compara quantos vizinhos
de cada \emph{tipo} a ocorrência tem nas duas redes; ele quase não depende do sorteio,
porque trocar uma cópia idêntica por outra não muda a composição.

% ------------------------------------------------------------
\subsection{Dominância lexical (P1 e P2)}\label{sec:met-dominancia}

A \textbf{dominância lexical} de uma ocorrência é a fração dos seus vizinhos que são
ocorrências do mesmo tipo,
\begin{equation}
    D_i=\frac{|\{\,j\in\N_i : t_j=t_i\,\}|}{k}.
\end{equation}
Na rede lexical, as outras ocorrências do tipo são as mais parecidas possíveis (cosseno 1), e
$D_i$ atinge o seu teto $\min(f_{t_i}-1,\,k)/k$: uma ocorrência de um tipo com 4 ocorrências
nunca passa de $0{,}3$. Para comparar tipos de frequências diferentes, usamos também a
dominância relativa ao teto, $\tilde D_i=D_i/\text{teto}$, que vale 1 na rede lexical. O nível
do acaso, se os vizinhos fossem sorteados, é $(f_{t_i}-1)/(n-1)$, no máximo
$49/13.605\approx0{,}004$.

% ------------------------------------------------------------
\subsection{Comunidades e concordância com rótulos (P3)}\label{sec:met-comunidades}

As comunidades são encontradas na rede por união com o algoritmo de
Leiden~\cite{traag2019}, que maximiza a modularidade~\cite{newman2006}. Ele roda 10 vezes
com sementes diferentes; fica a partição de maior modularidade, e a \textbf{estabilidade} é a
concordância média entre as 10 partições. O Louvain~\cite{blondel2008} serve de checagem.

Para saber o que as comunidades representam, comparamos a partição com rótulos conhecidos dos
vértices: o tipo de token, o tema, o artigo, o parágrafo, a sentença, o próximo token do texto,
o token que o modelo prevê e a faixa de posição no parágrafo. A posição entra como controle:
se ela explicasse as comunidades, o resultado seria um artefato da posição, e não contexto. A
comparação usa a \textbf{informação mútua normalizada} (NMI), que vai de 0 (partições
independentes) a 1 (iguais). O NMI cresce com o número de classes mesmo entre partições sem
relação nenhuma~\cite{vinh2010}, e os rótulos aqui vão de 4 classes (faixa de posição) a 2.780
(tipo de token). Por isso cada NMI é descontado de uma \textbf{linha de base}: o NMI médio
com o rótulo embaralhado 20 vezes, que preserva o número e o tamanho das classes. A medida
lida é NMI menos a linha de base. Reportamos também o índice de Rand
ajustado~\cite{hubert1985}, já corrigido pelo acaso.

A P3 é analisada \textbf{só nos \nVerticesNucleo{} vértices do núcleo}, em que os parágrafos
estão inteiros. As comunidades são as da rede inteira, restritas a esses vértices.

% ------------------------------------------------------------
\subsection{Separação das ocorrências de um mesmo tipo (P4)}\label{sec:met-p4}

Para cada tipo com pelo menos 10 ocorrências, separados em alvos, controles, multitema e
palavras funcionais (artigos, preposições, pronomes\ldots), medimos:
\begin{itemize}
    \item a \textbf{autossimilaridade}, o cosseno médio entre as ocorrências do tipo (vale 1 na
    camada lexical)~\cite{ethayarajh2019};
    \item a diferença \textbf{intra $-$ entre}: o cosseno médio entre ocorrências do mesmo tema
    menos o cosseno médio entre ocorrências de temas diferentes. Se for positiva, as
    ocorrências se agrupam por tema; se for zero, o tema não importa;
    \item o \textbf{número efetivo de comunidades} ocupadas pelas ocorrências, $e^{H}$, com $H$
    a entropia da distribuição delas pelas comunidades. Vale 1 quando todas estão na mesma
    comunidade;
    \item o \textbf{NMI entre comunidade e tema} dentro do tipo, com a sua linha de base: as
    ocorrências não só se espalham, mas se espalham \emph{segundo o tema}?
\end{itemize}
Nas palavras-alvo, as duas últimas medidas e a diferença intra $-$ entre são repetidas com o
\textbf{sentido anotado} no lugar do tema (Seção~\ref{sec:dados-sentidos}), em todas as
ocorrências e só nas que o contexto à esquerda decide.

% ------------------------------------------------------------
\subsection{Vizinhos no vocabulário}\label{sec:met-lente}

Como o Qwen3 usa a mesma matriz $E$ na entrada e na saída, podemos perguntar, para o vetor de
uma ocorrência em cada bloco, quais linhas de $E$ estão mais próximas dele. Ordenamos as
151.669 linhas pelo cosseno com o vetor e anotamos a posição do \emph{próprio} token e a do
\emph{próximo} token do texto nessa ordem. É uma versão simples da ``lente de
\emph{logits}''~\cite{nostalgebraist2020}. A pergunta é em que bloco a ocorrência deixa de se
parecer com o próprio token e passa a se parecer com o seguinte. Duas checagens validam o
cálculo: na camada lexical, o próprio token é sempre o primeiro; na saída normalizada do bloco
36, a linha de maior produto interno coincide com a previsão do próprio modelo em
\concordanciaPredicao{} das ocorrências.

% ============================================================
\section{Resultados}\label{sec:resultados}
% ============================================================

\previsto{Cerca de 4 páginas, uma subseção por pergunta, cada uma abrindo com a resposta em
uma frase. Escrita junto com a Seção 8 do relatório técnico, de onde vêm as figuras.}

% ------------------------------------------------------------
\subsection{Visão geral das redes}\label{sec:res-geral}

\previsto{O que muda nas métricas globais entre \cod{lex}, \cod{L01}, \cod{L18} e
\cod{L36}, e o que é esperado por construção.}

\tabelaopcional{metricas_globais}
    {Métricas estruturais das redes da configuração principal ($k=10$, união; média $\pm$
    desvio entre as 20 sementes na rede lexical, semente 0 nas camadas contextuais).}
    {tab:metricas-globais}
\figuraopcional{graus_camadas}
    {Distribuição do grau de entrada (CCDF, escala log-log) nas redes por ocorrência
    direcionadas, por representação. A linha vertical marca $k$, a média comum a todas.}
    {fig:graus-camadas}

% ------------------------------------------------------------
\subsection{P1: quanto e de que forma a organização muda}\label{sec:res-p1}

\previsto{Persistência da vizinhança e dominância lexical por camada, faixa de frequência e
categoria de token, contra o piso de ruído. Resposta à H1.}

\figuraopcional{p1_jaccard_faixas}
    {Distribuição de $J_i(\mathrm{lex},\ell)$ para $\ell\in\{1,18,36\}$ por faixa de frequência
    na amostra, com o piso de ruído $1-M^{\text{ruído}}$ de cada faixa.}
    {fig:p1-jaccard-faixas}
\figuraopcional{p1_dominancia}
    {Dominância lexical $D$ e dominância relativa ao teto $\tilde D$ por camada, faixa de
    frequência e categoria de token ($f_t\ge3$).}
    {fig:p1-dominancia}

% ------------------------------------------------------------
\subsection{P2: em que camadas a mudança se concentra}\label{sec:res-p2}

\previsto{Curvas camada a camada nas 36 camadas e em que transição a mudança se concentra,
por categoria de token. Resposta à H2.}

\figuraopcional{p2_curvas_camadas}
    {Curvas camada a camada: $J(\ell,\ell+1)$ entre blocos consecutivos e dominância
    $D(\ell)$, por categoria de token.}
    {fig:p2-curvas-camadas}

% ------------------------------------------------------------
\subsection{P3: comunidades e contexto}\label{sec:res-p3}

\previsto{Concordância das comunidades com tipo, tema, artigo, parágrafo, sentença e próximo
token, por camada, no núcleo. Resposta à H3.}

\figuraopcional{p3_nmi_camadas}
    {NMI menos a linha de base de permutação entre as comunidades de cada camada e cada rótulo
    (núcleo da amostra), na semente 0 do desempate. A estabilidade do Leiden (NMI médio entre
    as 10 execuções) é informada ao lado.}
    {fig:p3-nmi-camadas}

% ------------------------------------------------------------
\subsection{P4: separação das ocorrências de um mesmo token}\label{sec:res-p4}

\previsto{Separação intra-tipo por camada para alvos, controles e multitema; comparação entre
alvos e controles e separação pelo sentido anotado. Rede ego de \emph{banco}. Resposta à H4.}

\figuraopcional{p4_separacao}
    {Separação intra-tipo por camada: autossimilaridade, diferença entre cosseno intra e entre
    temas e NMI comunidade$\times$tema, para alvos, controles e multitema.}
    {fig:p4-separacao}
\figuraopcional{p4_ego_banco}
    {Rede ego de \tok{\esp banco} nas redes \cod{lex}, \cod{L01}, \cod{L18} e \cod{L36}, com
    as ocorrências coloridas pelo tema de sentido e posições fixas entre camadas.}
    {fig:p4-ego-banco}

% ============================================================
\section{Robustez}\label{sec:robustez}
% ============================================================

Os resultados só valem se não dependerem de uma escolha particular de construção. Refizemos as
redes mudando uma escolha de cada vez e repetimos as métricas estruturais e a troca de vizinhos
(Tabela~\ref{tab:robustez}). \textbf{Em nenhuma variação a conclusão qualitativa muda}: a
clusterização é maior no bloco 1 que nos blocos 18 e 36, as distâncias encurtam a partir do
bloco 18 e a troca de vizinhos em relação à rede lexical cresce do bloco 1 ao 36.

\tabelaopcional{robustez}
    {Métricas-resumo em cada variação, ao lado da configuração principal. $\bar C$:
    clusterização local média; $\bar\ell$: distância média; $Q$: modularidade do Leiden;
    $J$: Jaccard médio com a rede lexical de mesmo $k$ e mesma regra de empate. Cada linha
    muda uma escolha em relação à principal.}
    {tab:robustez}

\paragraph{Número de vizinhos.} Com $k=5$ ou $k=20$, os valores mudam na direção esperada
(mais vizinhos, distâncias menores: no bloco 18, 7,3 com $k=5$, 5,1 com $k=10$ e 4,0 com
$k=20$), mas a ordem entre as camadas se mantém. A dominância relativa ao teto no bloco 36
fica entre 0,37 e 0,43 nos três valores de $k$.

\paragraph{Rede mútua.} Na versão em que a aresta exige escolha dos dois lados, a rede se
fragmenta e as distâncias crescem (17,0, 9,1 e 10,4 nos blocos 1, 18 e 36), mas a ordem entre as
camadas é a mesma.

\paragraph{Empates.} O desempate pela posição dá praticamente os mesmos valores que o sorteio
(Jaccard de 0,359 contra 0,354 entre \cod{lex} e \cod{L01}). Aceitar todos os empatados
(variante b) aumenta um pouco o Jaccard (0,41 no bloco 1), porque a vizinhança lexical fica
maior e contém mais cópias do próprio token, mas mantém a ordem. Com uma tolerância de empate
cem vezes maior ($\varepsilon=10^{-4}$), as métricas não mudam na segunda casa decimal, e as
vizinhanças de 2 a 4\% das linhas das redes contextuais mudam. A
Figura~\ref{fig:robustez-empates} separa as variantes por faixa de frequência. Duas delas
completam a leitura da P1: o Jaccard por composição de tipos ($J^w$), que ignora o sorteio,
fica perto de 1 entre \cod{lex} e \cod{L01} para os tipos frequentes, o que confirma que o
bloco 1 mantém os vizinhos do mesmo token; e a variante (c), que ignora o próprio tipo e compara
só quais \emph{outros} tipos estão por perto, mostra que esses outros tipos já mudam no bloco 1
(Jaccard entre 0,3 e 0,4 em todas as faixas).

\figuraopcional{robustez_empates}
    {Jaccard médio entre a rede lexical e cada camada nas regras de empate (sorteio, posição,
    todos os empatados e tipos distintos) e Jaccard por composição de tipos ($J^w$), por faixa
    de frequência.}
    {fig:robustez-empates}

\paragraph{Normalização final e centragem.} A saída normalizada do bloco 36 (\cod{L36n}),
que é o que o modelo usa para prever o próximo token, dá redes quase iguais às da saída bruta
(clusterização 0,343 contra 0,353; Jaccard com a lexical 0,094 contra 0,096). Subtrair dos
vetores a média de todas as ocorrências, o que remove a direção comum (a anisotropia), também
muda pouco (no bloco 36, distância média de 5,72 contra 5,28). A direção comum não é o que
produz os resultados.

\paragraph{Comunidades.} Com resoluções 0,5 e 2 do Leiden, e com o Louvain, a modularidade das
redes \cod{L18} e \cod{L36} fica entre 0,73 e 0,76, perto da principal (0,76 e 0,77).

\paragraph{O que não foi testado.} As concordâncias da P3 e a separação da P4 foram medidas só
na configuração principal (e nos 36 blocos), e não nas variações acima. Uma rede construída só
com os vértices do núcleo, prevista no plano, não foi feita: ela exige refazer o k-NN, porque o
k-NN de um subconjunto não é um recorte do k-NN completo (quando os vizinhos de fora saem, cada
vértice precisa escolher outros).

% ============================================================
\section{Discussão}\label{sec:discussao}
% ============================================================

A Tabela~\ref{tab:respostas} junta as respostas às quatro perguntas.

\begin{table}[htbp]
\centering
\caption{As hipóteses e o que os resultados mostram.}\label{tab:respostas}
\small
\begin{tabularx}{\linewidth}{@{}>{\raggedright\arraybackslash}p{5.6cm} l X@{}}
\toprule
\textbf{Hipótese} & \textbf{Resultado} & \textbf{Evidência principal} \\
\midrule
H1. A dominância lexical cai, mas fica acima do acaso; surgem \emph{hubs}; as comunidades se
alinham menos ao token. & confirmada &
Dominância de 0,70 para 0,25 (acaso abaixo de 0,004); assimetria do grau de entrada de 0 para
4,9; concordância das comunidades com o token de 0,32 para 0,19. \\
\midrule
H2. Pedaços de palavra mudam antes das palavras inteiras. & em parte &
Continuações e números perdem a identidade até o bloco 4; palavras inteiras, aos poucos e
sobretudo depois do bloco 28; inícios de palavra quase não perdem. \\
\midrule
H3. As comunidades passam do token ao contexto e, no fim, ao próximo token. & confirmada, com
correção & O contexto tem máximo nos blocos 12 a 19 e cai depois; o próximo token sobe a partir
do bloco 21. \\
\midrule
H4. Palavras ambíguas se separam mais que as de controle. & confirmada &
NMI comunidade $\times$ tema de 0,41 contra 0,17 no bloco 36; separação maior pelo sentido do
que pelo tema em todas as 11 palavras. \\
\bottomrule
\end{tabularx}
\end{table}

\paragraph{Três regiões do modelo.} Juntas, as medidas descrevem o Qwen3-4B-Base em três
regiões. Nos primeiros blocos (até o 5, aproximadamente), o vetor de uma ocorrência ainda é
basicamente o seu token: a rede do bloco 1 é a rede lexical com os empates desfeitos, e a linha
mais próxima na matriz de \emph{embeddings} é a do próprio token. Nos blocos do meio (de 6 a
22, aproximadamente), o vetor deixa de se parecer com o próprio token e ainda não se parece com
o próximo; é nessa região que as comunidades seguem o artigo, o parágrafo e o tema, e que os
sentidos de uma palavra ambígua mais se afastam. No último terço, o vetor se aproxima do token
que vem depois, as comunidades passam a segui-lo, a anisotropia cresce e surgem \emph{hubs}
fortes. A transição entre as duas últimas regiões não é suave: nos blocos 22 a 25, a troca de
vizinhos entre blocos, a concordância com o contexto e a distância ao próximo token mudam ao
mesmo tempo. Essa divisão é coerente com análises do tipo ``lente de
\emph{logits}''~\cite{nostalgebraist2020}, em que as previsões do modelo só se formam nas
camadas finais.

\paragraph{O que as redes mostram além das médias de pares.} A autossimilaridade, a medida
clássica de quanto um token ficou contextual~\cite{ethayarajh2019}, cai até o meio do modelo e
volta a subir no fim, igual para palavras ambíguas e estáveis. Sozinha, ela sugeriria que o
último bloco é ``menos contextual'' que o do meio. As redes mostram outra coisa: a troca de
vizinhos continua até o último bloco, as comunidades mudam de natureza (do texto para o próximo
token), e a separação das palavras ambíguas pela rede cresce até o fim. A subida da
autossimilaridade acompanha a anisotropia, uma direção comum a todos os vetores, e não um
retorno à organização lexical.

\paragraph{A identidade do token é persistente.} A linha do próprio token em $E$ deixa de ser a
mais próxima já no bloco 5, mas as ocorrências de um mesmo token continuam vizinhas umas das
outras muito depois: a dominância lexical das palavras inteiras ainda é 0,64 no bloco 20. O
modelo mantém as ocorrências de um token juntas numa direção que não é mais o seu
\emph{embedding} de entrada.

\paragraph{O papel da fragmentação.} O tokenizer do Qwen3 divide muito o português, e isso
aparece nos resultados. Uma continuação de palavra (\tok{co} em ``ban\,co'') não significa nada
sozinha e perde a identidade logo nos primeiros blocos, quando o modelo junta os pedaços. Um
início de palavra mantém a identidade até o fim, provavelmente porque o próximo token, a
continuação da palavra, depende sobretudo dele. Com um tokenizer que preservasse mais palavras
inteiras, a curva média de reorganização provavelmente seria diferente.

\paragraph{Conclusão.} Redes k-NN sobre as mesmas ocorrências, uma por camada, permitem
acompanhar a reorganização das representações de um Transformer vértice a vértice e em várias
escalas. No Qwen3-4B-Base e em texto enciclopédico em português, essa reorganização é grande,
não monótona e dependente do tipo de token: a identidade lexical cede espaço aos poucos, o
contexto do texto domina o meio do modelo e o próximo token domina o fim, e as palavras
ambíguas separam os seus sentidos. A rede do vocabulário completa o quadro do lado da entrada:
ela é organizada por significado entre línguas, e não por escrita, e os seus \emph{hubs} são
tokens raros, provavelmente pouco treinados.

% ============================================================
\section{Limitações e trabalhos futuros}\label{sec:limitacoes}
% ============================================================

\previsto{Meia página, só as limitações que afetam a leitura dos resultados (fragmentação do
português pelo Qwen3, um único modelo, sentidos anotados por um modelo de linguagem, tamanho
da amostra). Relatório técnico, Seção 11.}

% ============================================================
\section*{Material complementar}
% ============================================================

\begin{itemize}
    \item Código, configurações e dados: \url{\repo}.
    \item Relatório técnico, com o método completo, as decisões, os problemas encontrados e
    todas as figuras e tabelas: \url{\tecnico}.
    \item Estudo de viabilidade feito antes do pipeline: \href{\viabilidade}{PDF no repositório}.
    \item Página interativa da coleta, do corpus e da amostra: \url{\coleta}.
\end{itemize}

\previsto{Referências bibliográficas, se a disciplina pedir.}


% ------------------------------------------------------------
\begin{thebibliography}{99}
\small

\bibitem{blondel2008}
V.~D. Blondel, J.-L. Guillaume, R.~Lambiotte e E.~Lefebvre.
Fast unfolding of communities in large networks.
\emph{Journal of Statistical Mechanics: Theory and Experiment}, P10008, 2008.

\bibitem{csardi2006}
G.~Csárdi e T.~Nepusz.
The igraph software package for complex network research.
\emph{InterJournal, Complex Systems}, 1695, 2006.

\bibitem{ethayarajh2019}
K.~Ethayarajh.
How contextual are contextualized word representations? Comparing the geometry of BERT,
ELMo, and GPT-2 embeddings.
Em \emph{Proceedings of EMNLP-IJCNLP}, 2019.

\bibitem{hubert1985}
L.~Hubert e P.~Arabie.
Comparing partitions.
\emph{Journal of Classification}, 2(1):193--218, 1985.

\bibitem{newman2006}
M.~E.~J. Newman.
Modularity and community structure in networks.
\emph{Proceedings of the National Academy of Sciences}, 103(23):8577--8582, 2006.

\bibitem{nostalgebraist2020}
nostalgebraist.
Interpreting GPT: the logit lens.
\emph{LessWrong}, 2020.

\bibitem{press2017}
O.~Press e L.~Wolf.
Using the output embedding to improve language models.
Em \emph{Proceedings of EACL}, 2017.

\bibitem{qwen3}
Qwen Team.
Qwen3 technical report.
\emph{arXiv:2505.09388}, 2025.

\bibitem{radovanovic2010}
M.~Radovanović, A.~Nanopoulos e M.~Ivanović.
Hubs in space: popular nearest neighbors in high-dimensional data.
\emph{Journal of Machine Learning Research}, 11:2487--2531, 2010.

\bibitem{sun2024}
M.~Sun, X.~Chen, J.~Z. Kolter e Z.~Liu.
Massive activations in large language models.
\emph{arXiv:2402.17762}, 2024.

\bibitem{traag2019}
V.~A. Traag, L.~Waltman e N.~J. van Eck.
From Louvain to Leiden: guaranteeing well-connected communities.
\emph{Scientific Reports}, 9:5233, 2019.

\bibitem{vinh2010}
N.~X. Vinh, J.~Epps e J.~Bailey.
Information theoretic measures for clusterings comparison: variants, properties,
normalization and correction for chance.
\emph{Journal of Machine Learning Research}, 11:2837--2854, 2010.

\end{thebibliography}

\end{document}
